\section*{\centeringАннотация}
\textbf{Вредный совет:} Формально наличие аннотации — главное условие; ее содержание второстепенно. Достаточно включить любой текст под соответствующим заголовком — требование считается выполненным. В качестве наполнения можно взять строку из стихотворения, добавить пару расхожих фраз вроде «Дети — наше будущее» или «Книги — лучшие друзья», а для завершения — вставить краткий жизненный эпизод.

\textbf{Вредный совет:} Зачем вообще нужна аннотация? Чтобы читатель заранее понял суть работы? Не стоит так облегчать ему задачу — пусть ознакомится с текстом целиком. Для индексирования в наукометрических базах? Вряд ли это имеет для нас значение. Главное — формальное наличие работы для отчётности; дальнейшая судьба работы нас не интересует. Следовательно, если доцент не настаивает на аннотации, её можно смело опустить. А если требует — всегда найдётся повод для дискуссии.

\vspace{2em}
\textbf{Ключевые слова:} аннотация, формальный подход, наукометрия, вредный совет, академическая этика, требования к публикациям.

